\documentclass[a4paper,11pt]{article}
\usepackage[utf8]{inputenc}
\usepackage[spanish]{babel}
\usepackage{amsmath, amsthm, amssymb, amsfonts}
\usepackage[dvipsnames]{xcolor}
\usepackage{tcolorbox}
\tcbuselibrary{skins, breakable, theorems}
\usepackage[margin=2.5cm]{geometry}

% Definición de colores
\definecolor{maincolor}{RGB}{41, 128, 185}
\definecolor{defcolor}{RGB}{39, 174, 96}
\definecolor{thmcolor}{RGB}{142, 68, 173}
\definecolor{notecolor}{RGB}{230, 126, 34}

% Configuración de las cajas
\newtcbtheorem[number within=section]{definicion}{}%
{colback=defcolor!5,colframe=defcolor!80!black,fonttitle=\bfseries, breakable}{def}

\newtcbtheorem[number within=section]{teorema}{Teorema}%
{colback=thmcolor!5,colframe=thmcolor!80!black,fonttitle=\bfseries, breakable}{thm}

\newtcolorbox{nota}[1][]{colback=notecolor!5,colframe=notecolor!80!black,fonttitle=\bfseries, breakable, title=#1}

\title{\textbf{\color{maincolor}Resumen I1 - Análisis Real}}
\author{Santiago Franulic}
\date{}

% Configuraciones extra
\renewcommand{\arraystretch}{1.4}
\usepackage{hyperref}
\hypersetup{
    colorlinks=true,
    linkcolor=maincolor,
    urlcolor=maincolor,
    citecolor=maincolor
}

\begin{document}
\maketitle
\tableofcontents
\newpage

\section{Espacios Metricos}

\subsection{Metricas}

\begin{definicion*}{Métrica $\rho$}{}
  Regla de asignación de distancias en un conjunto $X$
  \begin{itemize}
    \item Requisitos
    \[
    \begin{array}{lll}
      1. & (\rho(x,y) > 0) \Leftrightarrow (x \neq y) & \forall_{x \in X}\forall_{y \in X} \\
      2. & (\rho(x,y) = 0) \Leftrightarrow (x = y) & \forall_{x \in X}\forall_{y \in X} \\
      3. & \rho(x,y) = \rho(y,x) & \forall_{x \in X}\forall_{y \in X} \\
      4. & \rho(x,y) \leq \rho(x,z) + \rho(z,y) & \forall_{x \in X}\forall_{y \in X}\forall_{z \in X}
    \end{array}
    \]
    \item Ejemplos
    \begin{enumerate}
        \item Metrica discreta
        \begin{equation*}
            \rho(x,y) = \begin{cases}
                1 & \text{si } x \neq y \quad \forall_{x \in X}\forall_{y \in X} \\
                0 & \text{si } x = y \quad \forall_{x \in X}\forall_{y \in X}
            \end{cases}
        \end{equation*}
        \item Normas p
        \begin{equation*}
            \rho(x,y) = \left( \sum_{i=1}^{n} |x_i - y_i|^p \right)^{1/p}
        \end{equation*}
    \end{enumerate}
  \end{itemize}
\end{definicion*}

\begin{definicion*}{Espacio Métrico $(X,\rho)$}{}
  $ (X,\rho) \quad | \quad \rho \text{ es una métrica en } X $
\end{definicion*}

\begin{definicion*}{Subespacio Métrico $(Y,\rho_Y)$}{}
  $ (Y,\rho_Y) \quad | \quad \forall_{x \in Y} \forall_{y \in Y} \left( (Y \subseteq X) \wedge (\rho_Y(x,y) = \rho(x,y)) \right) $
\end{definicion*}

\begin{definicion*}{Convergencia de una Sucesión en un Espacio Métrico $(X,\rho)$}{}
  $\left( \lim_{n \to \infty} (x_n) = x \right) \Leftrightarrow \left(\lim_{n \to \infty} \rho(x_n,x) = 0 \right)$
\end{definicion*}

\begin{definicion*}{Métricas Equivalentes}{}
  $\left( \rho_1 \sim \rho_2 \right) \Leftrightarrow \begin{cases}
    (\rho_1 \text{ converge }) \Leftrightarrow (\rho_2 \text{ converge }) \\
    \lambda \rho_1 \leq \rho_2 \leq \mu \rho_1 \quad \forall_{x \in X}\forall_{y \in X}\exists_{\lambda \in \mathbb{R}^+} \exists_{\mu \in \mathbb{R}^+}
  \end{cases}$
\end{definicion*}

\newpage
\subsection{Normas}

\begin{definicion*}{Espacio vectorial}{}
  Conjunto de vectores con operaciones de suma y multiplicación por escalares
\end{definicion*}
  
\begin{definicion*}{Norma de un vector}{}
  Función que asigna una longitud o tamaño a cada vector
  \begin{itemize}
    \item Requisitos
    \[
    \begin{array}{lll}
      1. & \left( (\|x\| > 0) \Leftrightarrow (x \neq 0) \right) \wedge \left( (\|x\| = 0) \Leftrightarrow (x = 0) \right) & \forall_{x \in X} \\
      2. & \|\alpha x\| = |\alpha| \|x\| & \forall_{x \in X}\forall_{\alpha \in \mathbb{R}} \\
      3. & \|x+y\| \leq \|x\| + \|y\| & \forall_{x \in X}\forall_{y \in X}
    \end{array}
    \]
    \item Teorema de Cauchy-Schwarz
    \begin{equation*}
      |\langle x,y \rangle| \leq \|x\| \|y\| \quad \forall_{x \in X}\forall_{y \in X}
    \end{equation*}
  \end{itemize}
\end{definicion*}

\newpage
\subsection{Conjuntos Abiertos y Cerrados}

\begin{definicion*}{Bolas Abiertas y Cerradas}{}
  \[
  \begin{array}{lll}
    1. & \text{Bola abierta:} & B(x,r) = \{ y \in X \mid \rho(x,y) < r \} \\
    2. & \text{Bola cerrada:} & \bar{B}(x,r) = \{ y \in X \mid \rho(x,y) \leq r \} \\
  \end{array}
  \]
\end{definicion*}

\begin{definicion*}{Puntos interiores y limites}{} 
  \[
  \begin{array}{llll}
    1. & \text{Punto interior:} & x \in A &\mid \exists_{r > 0} (B(x,r) \cap A^c = \emptyset) \\
    2. & \text{Punto frontera:} & x &\mid \forall_{r > 0} ((B(x,r) \cap A \neq \emptyset) \wedge (B(x,r) \cap A^c \neq \emptyset)) \\
    3. & \text{Punto limite:} & x &\mid \forall_{r > 0} (B(x,r)\setminus\{x\} \cap A \neq \emptyset) \\
    4. & \text{Punto exterior:} & x &\mid \exists_{r > 0} (B(x,r) \cap A = \emptyset) \\
  \end{array}
  \]
\end{definicion*}

\begin{definicion*}{Conjuntos Abiertos y Cerrados}{}
  \[
  \begin{array}{lll}
    1. & \text{Conjunto abierto:} & A \mid \forall_{x \in A} (x \text{ es un punto interior}) \\
    2. & \text{Conjunto cerrado:} & A \mid \forall_{x} (x \text{ es un punto frontera} \Rightarrow x \in A)
  \end{array}
  \]
\end{definicion*}

\begin{definicion*}{Clausura e Interior}{}
  \[
  \begin{array}{llll}
    1. & \text{Clausura:} & \bar{A} &= \{ x \mid \forall_{r > 0} (B(x,r) \cap A \neq \emptyset) \} \\
    2. & \text{Interior:} & A^o &= \{ x \mid \exists_{r > 0} (B(x,r) \cap A^c = \emptyset) \} \\
  \end{array}
  \]
\end{definicion*}

\begin{definicion*}{Teoremas de Conjuntos Cerrados}{}
  \begin{itemize}
    \item $(F = \bar{F}) \Leftrightarrow \forall_{\text{sucesión convergente } (x_n) \subseteq F} \left( \lim_{n \to \infty} (x_n) \in F \right)$
    \item $(F = \bar{F}) \Leftrightarrow (F^c = (F^c)^o)$
    \item $(U = U_Y^o) \Leftrightarrow \exists_{G\subseteq X}(G = G_X^o \wedge G\cap Y = U) \qquad \forall_{U \subseteq Y \subseteq X} $
  \end{itemize}
  \begin{center}
  \begin{tabular}{l c c}
    \hline
    \textbf{Operación} & \textbf{Conjuntos Abiertos} & \textbf{Conjuntos Cerrados} \\
    \hline
    Unión finita & Abierto & Cerrado \\
    Intersección finita & Abierto & Cerrado \\
    Intersección infinita & No necesariamente abierto & Cerrado \\
    Unión infinita & Abierto & No necesariamente cerrado \\
    \hline
  \end{tabular}
  \end{center}
\end{definicion*}

\newpage

\section{Funciones Continuas en Espacios Métricos}

\subsection{Límites en Espacios Métricos}

\begin{definicion*}{Límite en un Espacios Métricos}{}
  \begin{equation*}
    \begin{array}{cc}
      \lim_{x \to x_0}(f(x)) = y_0 &\\
      \Leftrightarrow&\\
      \forall_x\left((0 < \rho(x, x_0) < \delta) \Rightarrow (\sigma(f(x),y_0) < \varepsilon) \right) &\qquad \forall_{\varepsilon > 0} \exists_{\delta > 0} 
    \end{array}
  \end{equation*}
  \begin{itemize}
    \item $x_0 \in A \vee x_0 \notin A$
    \item $(\exists$ Límite) $\Rightarrow$ (es único)
  \end{itemize}
\end{definicion*}

\begin{definicion*}{Teorema de Limites de Funciones con Sucesiones}{}
  \begin{equation*}
    \begin{array}{cc}
      \lim_{x \to x_0}(f(x)) = y_0 &\\
      \Leftrightarrow&\\
      (\lim_{n \to \infty} (x_n) = x_0) \Rightarrow (\lim_{n \to \infty} (f(x_n)) = y_0) & \qquad \forall_{(x_n) \subseteq A \setminus \{x_0\}}\\
    \end{array}
  \end{equation*}
\end{definicion*}

\newpage

\subsection{Funciones Continuas}

\begin{definicion*}{Continuidad de una Función}{}
  Función continua $\Leftrightarrow$ es continua en todo punto de su dominio
  \begin{equation*}
    \begin{array}{cc}
      f \text{ es continua en } x_0 &\\
      \Leftrightarrow&\\
      \lim_{x \to x_0}(f(x)) = f(x_0) &\\
    \end{array}
  \end{equation*}
\end{definicion*}

\begin{definicion*}{Álgebra de Funciones Continuas}{}
  Las siguientes operaciones mantienen la continuidad de las funciones
  \begin{center}
    \begin{tabular}{l l r}
      \hline
      \textbf{Operación} & \textbf{Espacios y Condiciones} & \textbf{Resultado} \\
      \hline
      Suma Ponderada & $f,g:X \to \mathbb{R}$, $\alpha, \beta \in \mathbb{R}$ & $\alpha f + \beta g:X \to \mathbb{R}$ \\
      Producto Puntual & $f,g:X \to \mathbb{R}$ & $f \cdot g:X \to \mathbb{R}$ \\
      Cociente & $f,g:X \to \mathbb{R}$, $g(x) \neq 0 \, \forall x \in X$ & $\frac{f}{g}:X \to \mathbb{R}$ \\
      Composición & $f:X \to Y$, $g:Y \to Z$ & $g \circ f:X \to Z$ \\
      \hline
    \end{tabular}
  \end{center}
\end{definicion*}

\begin{definicion*}{Homeomorfismos}{}
  $(f: X \to Y$ es homeomorfismo $) \Leftrightarrow$ ($f$ es biyectiva $\wedge$ $f$ es continua $\wedge$ $f^{-1}$ es continua).
\end{definicion*}

\begin{definicion*}{Isometría}{} 
  \begin{equation*}
  \begin{array}{c}
    (f: X \to Y \text{ es isometría }) \Leftrightarrow (f \text{ es biyectiva } \wedge \rho(x,y) = \sigma(f(x),f(y)))\\
    \therefore (f: X \to Y \text{ es isometría }) \Rightarrow (f \text{ es homeomorfismo })
  \end{array}
  \end{equation*}
\end{definicion*}

\begin{definicion*}{Funciones Lineales}{}
  \begin{equation*}
    f(\alpha_1 x_1 + \dots + \alpha_n x_n) = \alpha_1 f(x_1) + \dots + \alpha_n f(x_n) \qquad f: \mathbb{R}^n \to \mathbb{R}^m 
  \end{equation*}
  Ejemplo: Integral definida
\end{definicion*}

\end{document}
